\scanpage{116}
\mspage{57}
\margin{(174)\quad $\left\{\rule{0pt}{14em}\right.$}
\diagram{\begin{tikzcd}[ampersand replacement=\&, column sep=large,row sep=small]
D_{d,*} \arrow[r,"\sigma_d"] \arrow[d,equal,"\mathrm{def}"'] \& D_d \\
\displaystyle\coprod_{\substack{\deleted{\unsure{d'}}\\d'\in\mathbb N}}D_{d,d'}=\sigma^{-1}(D_d) \&
\end{tikzcd}}
qui fait de $D_{d,*}$ un esp. analytique sur $D_d$, \deleted{\unsure{prédisposé}} relativement ordonné (par la relation de factorisabilité
\diagram{\begin{tikzcd}[ampersand replacement=\&, column sep=small,row sep=small]
{} \& {} \arrow[dd,dotted] \\
\bullet \arrow[ru] \arrow[rd] \& \\
{} \& {}
\end{tikzcd}}
des flèches de source donnée) et loc. isomorphe à $\deleted{\unsure{E}}\unsure{\mathcal P}(I_d)^{\deleted{\unsure{0}}}_{D_d}$ \inserted{``pondéré'' par $d'$, avec sa relation d'ordre et la fonction ``cardinal''} (--- le faisceau constant d'ensembles ordonnés \unsure{simpliciaux} $\mathcal P(I_d)^0$ ---

Cela redonne donc un rev. $D'_d$ principal de groupe $\unsure{G_d|\mathfrak S_d}$, permettant de reconstituer les $D_{d,d'}$ --- mais pour reconstituer $D$, il faut de plus se donner \deleted{\unsure{les}} au moins les $D_{d,d'}\unsure{\text{---}|\longrightarrow}D_{d'}$, et (c'est plus sûr!) aussi
\[
D_{\unsure{d|d^\circ},d'}\mathbin{\times}_{D_{d'}}D_{d',d''}\unsure{\text{---}|\longrightarrow}D_{\unsure{d|d^\circ},d''}
\]
(la flèche composition), car il n'est pas clair, avec les seules hyp. faites jusqu'à présent, que les morphismes
\[
D_{d,d''}\unsure{\text{---}|\longrightarrow}D_d\times D_{d''}
\]
sont mono ?, i.e. que ``les flèches sont déterminées par sa source et par son but''. Mais finalement si, c'est une conséquence de c). Donc
\margin{115}
